\subsection{Invariant positive lengths and coherent word substitutions}

\begin{claim}{D.10}{Proved}{An invariant positive length forces passage}
\label{D:invariant-positive-length}
Use the concrete category, rich limit, group $G$, and directed
inclusion poset $\mathcal I$ of D.8. Let
$M:\mathcal C_\delta\to\mathrm{FreeMon}$ and
$S=\colim_{\mathcal I}M$. Suppose there is a monoid homomorphism
\[
 \ell:S\longrightarrow(\mathbb R_{\geq0},+),\qquad
 \ell(s)=0\Longleftrightarrow s=1,\qquad
 \ell(gs)=\ell(s)\quad(g\in G).
 \tag{D.10.1}
\]
Then $q_M$ is a weak equivalence, and, separately,
\[
                   T_i^{\mathcal C_\delta}(M_{\mathrm{ab}})=0
                   \qquad(i>0).
 \tag{D.10.2}
\]
Consequently, a based diagram $H:\mathcal C_\delta\to\Ab$ of
free abelian groups which has a natural real-valued linear
functional strictly positive on every chosen basis element and
has $T_i(H)\ne0$ for some $i>0$ admits no strict free-monoid lift
in those specified bases.
\end{claim}
\begin{proof}
By D.9, the filtered limit $S$ is cancellative and conical, and
every finite subset lies in a finite-rank free submonoid.
It is equidivisible as well: place the four factors of an equality
in a common free chart and apply Levi's lemma there.
We also use the stabilizer-generation identity verified
in R17 and D.4:
\[
                \langle G_E,G_F\rangle=G_{E\cap F}
 \tag{D.10.3}
\]
for finite closed substructures, including the empty structure.
Here $G_\varnothing=G$. Finite closed structures are rigid because
their orders are linear.

\emph{Fixed monoids are closed under taking factors.}
If $g(ab)=ab$, equidivisibility applied to $(ga)(gb)=ab$ makes
$ga$ and $a$ comparable for left divisibility. Their lengths
are equal by (D.10.1), so the intervening factor has length zero
and equals $1$. Thus $ga=a$, and cancellation gives $gb=b$.
In particular
\[
                    G_{ab}=G_a\cap G_b,
\]
where these are element stabilizers. Hence
$S_F=S^{G_F}$ is factor-closed in $S$: if $ab\in S_F$, then
$a,b\in S_F$. This is the only use of the length hypothesis.

\emph{Finite free charts and graph chains.}
Let $\mathcal J$ be the poset of finite-rank free submonoids
$H\subseteq S$. It is directed: the finite union of two free
bases lies in a finite free submonoid by D.9. Put
$H_F=H\cap S_F$. Factor-closedness shows that $H_F$ is the free
monoid on the subset
\[
             \operatorname{At}(H)_F
                  =\operatorname{At}(H)\cap S_F
\]
of the canonical atom basis of $H$. Indeed every letter of an
$H$-word lying in $S_F$ also lies in $S_F$. The union of the
$H_F$ is $S_F$. If $H\subseteq H'$, each letter of $H_F$
substitutes as a word on $\operatorname{At}(H')_F$.

Fix any monoid $C$ and a $G$-invariant homomorphism $\rho:S\to C$;
$G$ acts trivially on $C$. For $L\subseteq S$, let $D(L)$ be
the translation category with objects $C$ and arrows
\[
                    (s,c):c\longrightarrow\rho(s)c.
\]
For free $L$ this is the free category on its generator graph,
whose chain complex in degrees $1,0$ is
\[
 \Z[C\times\operatorname{At}(L)]\xrightarrow{\partial}\Z[C],
 \qquad \partial(c,s)=[\rho(s)c]-[c].
\]
Substitution sends an edge to the sum of the edges along the
substituted word. These chain maps compose exactly. Equivalently,
the map from the category's bar complex to its graph chain complex
sends a word-arrow to that path sum and all higher chains to zero;
it is natural for positive substitutions and is a homology
equivalence for a free category.

Define
\[
 \mathscr E=\colim_{H\in\mathcal J}
                  \Z[C\times\operatorname{At}(H)],\qquad
 \mathscr E_F=\colim_{H\in\mathcal J}
                  \Z[C\times\operatorname{At}(H)_F].
 \tag{D.10.4}
\]
At each chart the second module is a basis submodule of the
first. Filtered exactness therefore makes
$\mathscr E_F\to\mathscr E$ injective. This does not require
the transition maps of graph-chain modules to be injective;
they need not be.

Intrinsically, a word $s\in S$ and $c\in C$ determine a path
chain $D(s,c)$ in $\mathscr E$, independently of a chart after
passing to a common larger chart. These generators satisfy
\[
        D(st,c)=D(t,c)+D(s,\rho(t)c).
\]
The $G$ action sends $D(s,c)$ to $D(gs,c)$, and $G_F$ fixes
$\mathscr E_F$ pointwise. Extending closed embeddings to $G$
therefore defines a strict $A$-module $F\mapsto\mathscr E_F$,
independently of the extensions.

\emph{The support filtration.}
For finite closed $F,E_1,\ldots,E_r$ one has, as subgroups of
$\mathscr E$,
\[
 \mathscr E_F\cap\sum_{j=1}^r\mathscr E_{E_j}
                  =\sum_{j=1}^r\mathscr E_{F\cap E_j}.
 \tag{D.10.5}
\]
At a chart this is the distributivity identity for modules
spanned by subsets of one basis. An atom is fixed by both
$G_F$ and $G_E$ exactly when it is fixed by $G_{F\cap E}$,
by (D.10.3). Thus the basis intersections are the stated ones.
Filtered exactness preserves finite sums and intersections,
viewed as kernels inside the common ambient module, and proves
(D.10.5) after passage to (D.10.4).

Filter the module diagram by
\[
 \mathscr E_F^{\leq n}
       =\sum_{\substack{B\subseteq F\ \mathrm{closed}\\|B|\leq n}}
                        \mathscr E_B,
 \qquad
 V_B=\mathscr E_B\Big/
       \sum_{\substack{D\subsetneq B\ \mathrm{closed}}}\mathscr E_D.
\]
We set $\mathscr E_F^{\leq-1}=0$.
The empty support is allowed in these sums. They are finite at
each $F$. Distributivity gives the natural isomorphism
\[
 \mathscr E_F^{\leq n}/\mathscr E_F^{\leq n-1}
       \cong
       \bigoplus_{\substack{B\subseteq F\ \mathrm{closed}\\|B|=n}}V_B.
 \tag{D.10.6}
\]
For distinct supports of cardinality $n$, every intersection is
proper, so (D.10.5) identifies precisely the relations killed
by the lower-support quotient. Conversely every proper closed
subsupport already occurs in that quotient.

Strong homogeneity identifies $V_B$ along isomorphisms of
supports; $G_B$ acts trivially, and rigidity leaves no additional
finite automorphism action. Consequently, for $n>0$, (D.10.6)
is a direct sum, over isomorphism types of $n$-point structures
$B$, of the module diagrams
\[
                    \Z[A(B,-)]\otimes_\Z V_B.
\]
For $n=0$ it is the constant module $V_\varnothing$. All these
diagrams have vanishing higher derived colimits. The bar complex
of a representable summand is the chain complex of its element
category tensored with $V_B$; the initial object $(B,\id_B)$
gives a contraction for arbitrary abelian coefficients. The
constant case follows from contractibility of $BA$.
Induct on $n$ through the filtration's short exact sequences.
Since the filtration exhausts every $\mathscr E_F$ and the bar
construction commutes with filtered colimits, exactness of
filtered abelian colimits then gives
\[
                         T_i^A(\mathscr E_-) =0\qquad(i>0).
 \tag{D.10.7}
\]

\emph{Translation-category homology.}
The category $D(S_F)$ is the filtered colimit of the free
categories $D(H_F)$. Nerves commute with these colimits, and the
natural graph-chain calculation above yields
\[
 H_q(BD(S_F);\Z)=0\quad(q\geq2),
\]
and the exact sequence of $A$-modules
\[
 0\longrightarrow H_1(BD(S_-);\Z)
 \longrightarrow\mathscr E_-
 \longrightarrow\Z[C]
 \longrightarrow\Z[\pi_0D(S_-)]\longrightarrow0.
 \tag{D.10.8}
\]
The constant term $\Z[C]$ and the final free Set-action module
are acyclic for higher colimits, by contractibility of $BA$ and
totality of $A$. Split (D.10.8) at the image of its boundary.
The two long exact sequences, together with (D.10.7), give
\[
              T_i^A\bigl(H_1(BD(S_-);\Z)\bigr)=0\qquad(i>0).
\]
In the Grothendieck homology spectral sequence for $D(S_-)$,
the rows $q\geq2$ vanish, the row $q=0$ has no terms in
positive column degree, and the same holds for row $q=1$.
Therefore
\[
                 H_j\left(B\int_A D(S_F);\Z\right)=0
                 \qquad(j\geq2).
 \tag{D.10.9}
\]

\emph{The comparison comma is contractible.}
Take $C=\colim_A^{\mathrm{Mon}}T_S$ and its invariant universal
map $\rho:S\to C$. The categorical colimit of $D(S_F)$ is the
regular translation category $C//C$. Indeed its generators are
stage elements acting at each object $c\in C$, and its transition
and multiplication relations are precisely the monoid-colimit
relations in each such context. Its arrows are thus $c\to mc$
for $m\in C$. The object $1$ is initial.

The fixed-object transport argument of D.6 applies to this
diagram: $BA$ is simply connected, so groupoid completion of its
Grothendieck construction agrees with groupoid completion of its
categorical colimit. It follows that $B\int_A D(S_F)$ is
connected and simply connected. Together with (D.10.9), Hurewicz
and Whitehead make it weakly contractible. This construction is
the sole comma category of $q_{T_S}$, so Theorem A proves that
$q_{T_S}$ is a weak equivalence. D.8 transports this conclusion
to $q_M$.

\emph{The separate abelianization assertion.}
Use the same chain construction with $C$ the one-element monoid.
Then $\mathscr E_F=(S_F^{\mathrm{gp}})_{\mathrm{ab}}$, the
universal abelian group of $S_F$, and its boundary is zero.
Equation (D.10.7) makes this abelianization diagram acyclic for
higher colimits. Each $BS_F$ has no homology above degree one,
as it is the filtered colimit of the free-monoid charts. Its
homotopy-colimit spectral sequence therefore gives
$H_j(B\int_A T_S;\Z)=0$ for $j\geq2$. D.8 identifies this
with the corresponding homology of $B\int_A M$, and D.1 gives
(D.10.2).

Finally, a strict word lift of a based diagram $H$ as in the
statement would inherit a positive length from the specified
natural functional. It remains strictly positive in the filtered
limit: a represented nonidentity word has positive length,
preserved by every transition. Naturality also makes the induced
length $G$-invariant. Equation (D.10.2) would contradict
$T_i(H)\ne0$. This proves the precise based nonlift assertion.
\end{proof}

The positive stabilization in the working note
\texttt{D\_positive\_stabilization.md} illustrates the last
assertion. Its specified bases carry the natural positive
functional $L_X=\lambda^{-N_X}\ell B_X$, while its first derived
colimit is $\Z$. Thus its positive matrices cannot be replaced
by strictly composable words in those bases. This does not rule
out different positive bases on the same abstract abelian diagram,
and (D.10.1) is an extra hypothesis on a general monoid diagram.
The unrestricted $\mathcal C_\delta$ comparison remains open here.
